% Section 4: Maximum Likelihood Estimation
\section{Maximum Likelihood Estimation}
\label{sec:estimation}

Given $n$ independent observations $t_1,\ldots,t_n$ from the parallel
system $T = \max(T_1,\ldots,T_m)$ with $T_j \sim \mathrm{Exp}(\lambda_j)$,
the observed-data log-likelihood under Scheme~0 is
\eqref{eq:loglik-exact}, restated here for reference:
\[
  \ell(\blam)
  = \sum_{i=1}^{n} \ln f_T(t_i;\blam)
  = \sum_{i=1}^{n} \ln \Biggl[\,\sum_{j=1}^{m} w_j(t_i;\blam)\Biggr].
\]
\Cref{thm:identifiability} guarantees that the maximizer is unique
(up to permutation of equal rates), and \Cref{thm:fisher} ensures
that the Fisher information is positive definite at distinct rates, so the
MLE is well-defined.

We present two approaches to computing the MLE: direct numerical
maximization (\S\ref{sub:direct-mle}) and an EM algorithm with
closed-form steps (\S\ref{sub:em-algorithm}).  Both converge to the
same estimate; the EM is particularly attractive because it avoids
gradient computation and is guaranteed to increase the likelihood at
every iteration.

% ===================================================================
\subsection{Direct maximization}
\label{sub:direct-mle}
% ===================================================================

The log-likelihood is a smooth function of
$\blam \in \Rset_{>0}^m$, and its gradient and Hessian are available
analytically through the inclusion-exclusion expansion of
\Cref{lem:ie}.  The score vector has components
\begin{equation}
  \label{eq:score}
  \frac{\partial \ell}{\partial \lambda_j}
  = \sum_{i=1}^{n}
    \frac{1}{f_T(t_i)}
    \Biggl[
      \frac{\partial w_j(t_i)}{\partial \lambda_j}
      + \sum_{k \neq j} \frac{\partial w_k(t_i)}{\partial \lambda_j}
    \Biggr],
\end{equation}
where the partial derivatives~\eqref{eq:dwj-dlamj}--\eqref{eq:dwk-dlamj}
are finite sums of exponentials, so the score and observed Fisher information
$I_{\mathrm{obs}}(\hat\blam) = -\nabla^2 \ell(\hat\blam)$
are analytically tractable.  Standard errors are obtained as
$\widehat{\mathrm{se}}(\hat\lambda_j)
= \bigl[I_{\mathrm{obs}}^{-1}(\hat\blam)\bigr]_{jj}^{1/2}$.

In practice, we maximize~$\ell(\blam)$ using the L-BFGS-B algorithm
with box constraints $\lambda_j > 0$ for $j = 1,\ldots,m$.  As a
fallback for robustness, we re-optimize with the
Nelder--Mead simplex method on the reparametrized scale
$\log\lambda_j \in \Rset$, which enforces positivity without box
constraints; the solution yielding the higher log-likelihood is
retained.

\paragraph{Starting values.}
For a parallel system of $m$ i.i.d.\ $\mathrm{Exp}(\lambda)$ components,
$\E[T] = H_m / \lambda$ where $H_m = \sum_{k=1}^{m} 1/k$ is the $m$-th
harmonic number.  Writing $\hat\lambda_0 = H_m / \bar{t}$ for the
method-of-moments estimate (where $\bar{t}$ is the sample mean of the
observed system lifetimes), we initialize with a spread:
\begin{equation}\label{eq:spread-init}
  \lambda_j^{(0)} = \hat\lambda_0 \cdot
  \Bigl(\tfrac{1}{2} + \tfrac{j-1}{m-1}\Bigr), \qquad j = 1,\ldots,m.
\end{equation}
This places the initial rates in the correct order of magnitude while
breaking the permutation symmetry of the Scheme~0 log-likelihood.
A uniform initialization $\lambda_j^{(0)} = \hat\lambda_0$ for all~$j$
sits at a saddle point of the likelihood surface---where all rates are
equal---and gradient-based optimizers fail to escape.

% ===================================================================
\subsection{EM algorithm}
\label{sub:em-algorithm}
% ===================================================================

The censoring equivalence (\Cref{prop:censoring-equiv})
suggests a natural data-augmentation scheme: augment each observed
system lifetime~$T$ with the latent identity $J = \argmax_j T_j$ and
the unobserved component lifetimes $T_1,\ldots,T_m$.  This leads to an
EM algorithm in which both the E-step and the M-step admit closed-form
expressions.

\begin{proposition}[EM algorithm for exponential parallel systems]
  \label{prop:em-algorithm}
  Let $t_1,\ldots,t_n$ be $n$ i.i.d.\ observations from
  $T = \max(T_1,\ldots,T_m)$ with $T_j \sim \mathrm{Exp}(\lambda_j)$,
  $\lambda_j > 0$.  Define the complete data as
  $(T_1,\ldots,T_m,\, J)$ where $J = \argmax_j T_j$.

  Given current parameter estimates $\blam^{(s)}$, the EM update
  $\blam^{(s)} \mapsto \blam^{(s+1)}$ proceeds as follows.

  \medskip\noindent
  \textbf{E-step.}\enspace
  For each observation $t_i$ and each component $j = 1,\ldots,m$,
  compute:
  \begin{enumerate}
    \item[\textup{(i)}]
      The posterior probability that component $j$ was the last to fail:
      \begin{equation}
        \label{eq:em-posterior}
        p_{ij}
        \;\coloneqq\; P\bigl(J = j \mid T = t_i;\, \blam^{(s)}\bigr)
        = \frac{w_j\bigl(t_i;\, \blam^{(s)}\bigr)}
              {f_T\bigl(t_i;\, \blam^{(s)}\bigr)}\,.
      \end{equation}

    \item[\textup{(ii)}]
      The conditional mean of $T_j$ given that component $j$ was the
      last to fail:
      \begin{equation}
        \label{eq:em-exact}
        \E\bigl[T_j \mid T = t_i,\, J = j;\, \blam^{(s)}\bigr] = t_i\,.
      \end{equation}

    \item[\textup{(iii)}]
      The conditional mean of $T_j$ given that component $j$ was
      \emph{not} the last to fail:
      \begin{equation}
        \label{eq:em-truncated}
        \E\bigl[T_j \mid T = t_i,\, J \neq j;\, \blam^{(s)}\bigr]
        = \E\bigl[T_j \mid T_j < t_i;\, \lambda_j^{(s)}\bigr]
        = \frac{1}{\lambda_j^{(s)}}
          - \frac{t_i}{e^{\lambda_j^{(s)}\, t_i} - 1}\,.
      \end{equation}
  \end{enumerate}
  The overall conditional expectation of $T_j$ given observation $t_i$
  is then
  \begin{equation}
    \label{eq:em-etj}
    \hat\tau_{ij}
    \;\coloneqq\; \E\bigl[T_j \mid T = t_i;\, \blam^{(s)}\bigr]
    = p_{ij}\cdot t_i
      + (1 - p_{ij})\biggl[
        \frac{1}{\lambda_j^{(s)}}
        - \frac{t_i}{e^{\lambda_j^{(s)}\, t_i} - 1}
      \biggr].
  \end{equation}

  \medskip\noindent
  \textbf{M-step.}\enspace
  Update each rate as
  \begin{equation}
    \label{eq:em-mstep}
    \lambda_j^{(s+1)}
    = \frac{n}{\displaystyle\sum_{i=1}^{n} \hat\tau_{ij}}\,,
    \qquad j = 1,\ldots,m.
  \end{equation}
\end{proposition}

\begin{proof}
We verify each component of the E-step and derive the M-step.

\smallskip\noindent
\emph{Part~\textup{(i)}.}\enspace
By \Cref{lem:density-decomp},
$w_j(t;\blam)/f_T(t;\blam) = P(J = j \mid T = t;\blam)$.
Evaluating at $\blam = \blam^{(s)}$ gives~\eqref{eq:em-posterior}.

\smallskip\noindent
\emph{Part~\textup{(ii)}.}\enspace
Conditioning on $J = j$ means that component~$j$ achieved the maximum:
$T_j = T = t_i$.  Hence
$\E[T_j \mid T = t_i,\, J = j] = t_i$.

\smallskip\noindent
\emph{Part~\textup{(iii)}.}\enspace
Conditioning on $J \neq j$ means $T_j < T = t_i$, so $T_j$ follows
an $\mathrm{Exp}(\lambda_j)$ distribution truncated above at~$t_i$.
Its conditional density on $(0,\, t_i)$ is
$g(u) = \lambda_j\, e^{-\lambda_j u}\big/\bigl(1 - e^{-\lambda_j t_i}\bigr)$.
The conditional mean is
\begin{align*}
  \E[T_j \mid T_j < t_i]
  &= \frac{\displaystyle\int_0^{t_i} u\,\lambda_j\, e^{-\lambda_j u}\,du}
          {1 - e^{-\lambda_j t_i}}
  = \frac{1/\lambda_j - (1/\lambda_j + t_i)\,e^{-\lambda_j t_i}}
          {1 - e^{-\lambda_j t_i}}
  = \frac{1}{\lambda_j}
    - \frac{t_i}{e^{\lambda_j t_i} - 1}\,,
\end{align*}
where the numerator follows from integration by parts and the final
form is obtained by dividing numerator and denominator by
$e^{-\lambda_j t_i}$.

\smallskip\noindent
\emph{M-step derivation.}\enspace
The complete-data log-likelihood for a single system~$i$ with augmented
data $(T_{i1},\ldots,T_{im},\, J_i)$ is
\[
  \ell_c^{(i)}(\blam)
  = \sum_{j=1}^{m}\bigl(\ln \lambda_j - \lambda_j\, T_{ij}\bigr)
    + \text{terms not depending on~$\blam$},
\]
since each $T_{ij} \sim \mathrm{Exp}(\lambda_j)$ independently.  The
complete-data sufficient statistic for~$\lambda_j$ is
$\sum_{i=1}^n T_{ij}$, and the complete-data MLE is
$\hat\lambda_j = n \big/ \sum_{i} T_{ij}$.  Taking the conditional
expectation given the observed data and current parameters yields
$\E\bigl[\sum_i T_{ij} \mid t_1,\ldots,t_n;\, \blam^{(s)}\bigr]
= \sum_i \hat\tau_{ij}$.  The M-step maximizes the expected
complete-data log-likelihood, giving~\eqref{eq:em-mstep}.

\smallskip\noindent
\emph{Monotonicity.}\enspace
By the general theory of \citet{DempsterLairdRubin1977},
$\ell\bigl(\blam^{(s+1)}\bigr) \geq \ell\bigl(\blam^{(s)}\bigr)$ at
every iteration, with equality if and only if $\blam^{(s)}$ is a fixed
point of the EM map.  The complete-data model belongs to the exponential
family, so each M-step has a unique closed-form maximum.
\end{proof}

\begin{remark}[Closed-form character]
  \label{rem:em-closed-form}
  Both the E-step and M-step require only elementary operations:
  the E-step evaluates $w_j$ via the inclusion-exclusion expansion
  ($2^{m-1}$ terms per component) and the truncated exponential
  mean~\eqref{eq:em-truncated}; the M-step computes a reciprocal
  of a weighted average.  No numerical integration or Monte Carlo
  sampling is needed at any step.  This contrasts with the stochastic
  EM of \citet{NgNavarroBalakrishnan2017}, which requires simulation
  in the E-step for general system signatures, and with the Bayesian
  approach of \citet{BheringPereiraPolpo2014}, which uses MCMC.
\end{remark}

\begin{algorithm}[t]
  \caption{EM algorithm for exponential parallel system MLE}
  \label{alg:em}
  \begin{algorithmic}[1]
    \Require Observations $t_1,\ldots,t_n > 0$;\;
             number of components~$m$;\;
             tolerance $\varepsilon > 0$
    \Ensure  MLE $\hat\blam = (\hat\lambda_1,\ldots,\hat\lambda_m)$
    \State $\hat\lambda_0 \gets H_m / \bar{t}$
      \Comment{$H_m = \sum_{k=1}^m 1/k$,\; $\bar{t} = n^{-1}\sum_i t_i$}
    \State Initialize $\lambda_j^{(0)} \gets \hat\lambda_0 \cdot
      \bigl(\tfrac{1}{2} + \tfrac{j-1}{m-1}\bigr)$ for $j = 1,\ldots,m$
      \Comment{spread init~\eqref{eq:spread-init}}
    \State $s \gets 0$
    \Repeat
      \For{$i = 1,\ldots,n$;\; $j = 1,\ldots,m$}
        \Comment{\textbf{E-step}}
        \State $p_{ij} \gets w_j(t_i;\blam^{(s)}) \,\big/\, f_T(t_i;\blam^{(s)})$
        \State $\mu_{ij} \gets 1/\lambda_j^{(s)} - t_i\big/(e^{\lambda_j^{(s)} t_i} - 1)$
          \Comment{truncated mean}
        \State $\hat\tau_{ij} \gets p_{ij}\cdot t_i + (1 - p_{ij})\cdot \mu_{ij}$
      \EndFor
      \For{$j = 1,\ldots,m$}
        \Comment{\textbf{M-step}}
        \State $\lambda_j^{(s+1)} \gets n \,\big/\, \sum_{i=1}^n \hat\tau_{ij}$
      \EndFor
      \State $s \gets s + 1$
    \Until{$\bigl|\ell(\blam^{(s)}) - \ell(\blam^{(s-1)})\bigr| \big/ \bigl|\ell(\blam^{(s-1)})\bigr| < \varepsilon$}
    \State \Return $\hat\blam \gets \blam^{(s)}$
  \end{algorithmic}
\end{algorithm}

% ===================================================================
\subsection{Practical considerations}
\label{sub:practical}
% ===================================================================

\paragraph{Convergence criterion.}
We terminate the EM when the relative change in log-likelihood falls
below a threshold~$\varepsilon$:
$|\ell(\blam^{(s)}) - \ell(\blam^{(s-1)})| / |\ell(\blam^{(s-1)})| < \varepsilon$.
In our simulations (\Cref{sec:simulations}), $\varepsilon = 10^{-8}$
was used, and convergence typically required 50--200 iterations.

\paragraph{Agreement of methods.}
Both the direct maximization and the EM
algorithm target the same log-likelihood~\eqref{eq:loglik-exact}.  Since
\Cref{thm:identifiability} guarantees a unique global maximum
(up to label permutation), they converge to the same point estimate.
The direct method supplies the observed Fisher information at the MLE
as a byproduct; when using the EM, standard errors are obtained by
evaluating the Hessian of~$\ell$ at the converged estimate~$\hat\blam$.

\paragraph{Computational cost.}
Each evaluation of $f_T(t;\blam)$ requires
the inclusion-exclusion expansion for each of $m$ components, involving
$2^{m-1}$ terms per component.  Over $n$ observations, a single
likelihood evaluation costs $O(n \cdot m \cdot 2^{m-1})$.  The
exponential scaling in~$m$ is the binding constraint: for the
exponential model, systems with $m \leq 15$--$20$ are computationally
feasible on standard hardware.

% ===================================================================
\subsection{Extension to censored system data}
\label{sub:censored-extension}
% ===================================================================

In practice, the system lifetime~$T$ may itself be subject to censoring.
The inclusion-exclusion machinery of \Cref{sec:likelihood}
accommodates all standard censoring types in closed form.

If system~$i$ is known only to survive past time~$\tau_i$
(i.e., $T_i > \tau_i$), its likelihood contribution is the system
survival function
$S_T(\tau_i;\blam) = 1 - F_T(\tau_i;\blam)$,
available in closed form via \Cref{lem:ie}.
If the system failure is known to lie in an interval $(a_i,\, b_i]$,
the likelihood contribution is
$F_T(b_i;\blam) - F_T(a_i;\blam)$,
a difference of two inclusion-exclusion expansions.

For a dataset containing exact, right-censored, and interval-censored
observations, the mixed log-likelihood~\eqref{eq:loglik}
applies directly, with every term in closed form.  Both the direct
maximization and EM approaches extend to this setting without
modification.  For the EM, the E-step conditions on the appropriate
censoring information: exact observations contribute as in
\Cref{prop:em-algorithm}; censored observations contribute the
conditional expectation of the component lifetimes given the
censoring constraint, which involves only exponential and truncated
exponential calculations.
